%==============================================================================
% discussion.tex — Discussion for PP-HFM
% Self-contained \section{}; \input this into main.tex (never compile alone).
%
% STATUS: PARTIAL. Only the Limitations subsection is written. The interpretive
% subsections are deliberately withheld until Results is complete, because
% three of their arguments are still undetermined:
%   - whether the cost of the privacy unit scales with records per individual
%     (records-per-individual experiment, running)
%   - whether every component earns its place (component ablation A1--A4,
%     never run)
%   - whether the fusion gain survives more seeds  [RESOLVED 2026-08-16: it
%     does, at n=5, but shrank from +0.0411/+0.1062 to +0.0370/+0.0671 and
%     from 14.1x to 5.9x seed-sd. Quote the five-seed figures only.]
% Writing the interpretation against a half-populated Results guarantees a
% rewrite. Every limitation below is settled and will not move.
%==============================================================================

\section{Discussion}
\label{sec:discussion}

\subsection{Limitations}
\label{sec:discussion:limits}

The claims of this paper are bounded in ways that are worth stating precisely,
because several of them constrain what may be generalised from the results
rather than merely qualifying them.

\textbf{The privacy-unit results are single-domain, for a structural reason.}
The framework is evaluated on two corpora, but the measurements concerning the
unit of the guarantee are reported only on the clinical one. This is not an
oversight of scope. In the wearable corpus each client corresponds to exactly
one participant, so client-level, user-level and individual-level guarantees
coincide and the divergence the mechanism addresses cannot arise there. A
second corpus exhibiting the same many-individuals-per-client structure --- a
second clinical cohort, or an institutional deployment in another domain ---
would be required to establish that the measured cost transfers, and its
absence is the most predictable objection to this work.

\textbf{The fusion result is single-domain, and the second corpus cannot
corroborate it either.} The gain from combining the modalities is established
on the clinical cohort, where it reaches $0.0370$ in area under the ROC curve
over the better single modality across five seeds, or $5.9$ times the seed
standard deviation. It is not established on the wearable cohort, and the
reason is a property of that task rather than of the framework: its better
single modality already attains $0.9976$, leaving a gain of $0.0004$ and no
headroom in which fusion could demonstrate anything. This is worth reporting
rather than omitting, because it qualifies the premise the framework is built
on. Multimodal federated learning is not universally motivated; on tasks a
single modality already solves, the second stream adds cost without adding
information, and the case for fusion must be made per task rather than assumed
from the presence of two sensors. Establishing that the measured gain transfers
would require a second corpus that is both genuinely paired at source and
difficult enough to leave room for the second modality to matter --- the same
requirement, and the same absence, as for the privacy-unit results above.

\textbf{Reported utilities come from the oracle encoder configuration.} As
Section~\ref{sec:exp:impl} records, the encoders behind the reported results
were pretrained on the pooled client training data, which does not carry a
valid end-to-end guarantee. The deployable configuration costs $0.0322$ in
area under the ROC curve, and every absolute utility in this paper should be
read as an upper bound on what it attains. The differences between arms, on
which the claims rest, are unaffected because every arm shares the same
encoders; but a reader comparing these figures against a system that respected
disjointness throughout would not be comparing like with like.

\textbf{The membership-inference null bounds a large effect only.} The audit
establishes that the difference in leakage between a record-level and an
individual-level guarantee is not of the magnitude hypothesised. It does not
establish that the difference is zero. At five paired seeds the half-width on
the strongest attack is $\pm 0.060$ in area under the curve, three of the four
attacks point in the predicted direction, and the strongest of them reaches
$p = 0.090$; an effect below roughly six points would be invisible to this
design. The result should be cited as a refutation of the predicted effect,
never as evidence that the two guarantees are equivalent.

\textbf{The audit rests on sixteen shadow models.} Per-individual in- and
out-distributions are estimated from roughly eight observations each, which is
few for a variance estimate and is the reason a pooled-variance variant is
reported alongside the per-individual one. A larger shadow population would
tighten those fits and is the first thing to increase if the audit were to be
repeated. Relatedly, with $600$ non-members the receiver operating
characteristic cannot resolve a $0.1\,\%$ false-positive rate, so the
low-false-positive corner --- which is where a confident disclosure would
appear --- is characterised only at $1\,\%$.

\textbf{One objective is inconclusive and must remain so.} On the wearable
deployment the accuracy cost under the budget is $1.967$ pp with an interval of
$[1.760, 2.173]$, which spans the $2$ pp requirement. The seed count of fifteen
was fixed in advance, and extending it until the interval fell on the
favourable side of the threshold would convert a pre-registered test into a
search. The objective is therefore reported as undetermined rather than as
either met or failed.

\textbf{Several supporting measurements rest on few seeds.} The privacy-budget
sweep of Fig.~\ref{fig:epsilon} is single-seed at each point, so the claim it
supports is the monotonicity of the trend rather than any individual value.
The aggregation-rule comparison that led to a component being dropped was
likewise decided on one seed; the effect was small enough that additional
seeds were judged unlikely to change the decision, but that judgement is not a
measurement.

\textbf{The network model simulates the network only.} Communication delay,
delivery reliability and traffic volume are measured on a packet-level
simulation of the topology, but the compute terms that enter the end-to-end
round latency --- sensing, preprocessing, inference and aggregation --- are
modelled analytically and are common to every topology compared. The
comparison between topologies is therefore sound, since the modelled terms
cancel, but the absolute round latency is not a measurement of a deployed
system and should not be read as one.

\textbf{Interpretability evidence is associational.} The attribution analysis
establishes that the model attends substantially to indicators of whether a
measurement was taken, and that this reliance increases under the privacy
mechanism. It does not establish that the model would lose accuracy without
those indicators, which only their removal could show. The claim is therefore
stated as evidence of reliance, not as a causal contribution to performance.

%==============================================================================
% PENDING SUBSECTIONS — write once Results is complete.
%
%   Interpretation of the privacy-unit cost
%     needs : records-per-individual scaling law (running)
%     turns on whether the cost is a constant of this cohort or a function of
%     records per individual; the second is a far stronger statement
%
%   Necessity of the components
%     needs : component ablation A1--A4 (never run)
%
%   Position against prior art
%     needs : the above, plus the fusion ablation at more than two seeds
%
%   Practical implications for deployment
%==============================================================================
